% Preamble additions for the pandoc/xelatex manual build (manual.pdf).
% Included via --include-in-header from the PDF build in .github/workflows/
% docs.yml and docs-publish.yml; the PDF manual only, never the ePub.

% --- Page -------------------------------------------------------------------
% pandoc's LaTeX template sets neither a paper size nor a geometry, so the
% manual inherited the article class's US-letter defaults: a 121 mm text column
% on 216 mm paper, with ~47 mm of white down each side and ~42 mm top and
% bottom. That is a page of margin with a manual in the middle of it.
%
% A4 with the margins a printed manual normally carries: 25 mm on three sides,
% a deeper 30 mm foot so the folio sits in air rather than under the last line.
% geometry loads AFTER the class options, so it — not the letterpaper default —
% decides the PDF's page size.
\usepackage[a4paper,left=25mm,right=25mm,top=25mm,bottom=30mm,footskip=14mm]{geometry}

% --- Code blocks ------------------------------------------------------------
% fancyvrb, which pandoc's default template uses, prints an over-long code line
% straight off the right edge of the page — silently, without even an overfull
% warning. fvextra teaches both verbatim environments to wrap: at a space where
% it can, mid-token when it must, with a hooked continuation marker so a reader
% can see that the line was wrapped rather than written that way.
%
% `breaknonspaceingroup` is the option that matters for a syntax-highlighted
% block: without it fvextra can only break BETWEEN pandoc's highlight macros,
% so a single long token inside one of them (a URL, a sha256, a one-line JSON
% array) still runs off the page.
\usepackage{fvextra}
\DefineVerbatimEnvironment{Highlighting}{Verbatim}{%
  commandchars=\\\{\},breaklines,breakanywhere,breaknonspaceingroup}
\RecustomVerbatimEnvironment{verbatim}{Verbatim}{%
  breaklines,breakanywhere,breaknonspaceingroup}

% --- Tables -----------------------------------------------------------------
% pandoc gives every column an equal share of the line, so the reference tables
% with four columns (agent, model, tools, description) get roughly 4 cm each —
% narrow enough that even a broken-up token lands past the column edge. A step
% down in size buys about a fifth more room per column and is the usual
% treatment for a dense reference table.
\AtBeginEnvironment{longtable}{\small}

% --- Characters the bundled fonts have no glyph for -------------------------
% xelatex drops an unmapped character silently, so "≥90% new-code coverage"
% prints as "90% new-code coverage" and "plugin.json ↔ marketplace.json" loses
% its arrow. Map the ones the docs actually use onto glyphs Latin Modern and
% amssymb do have. `newunicodechar` is not in the pandoc image's TeX Live set,
% so declare the active characters directly.
%
% The docs workflows fail the build on any REMAINING dropped character,
% so a newly introduced one lands here rather than vanishing from the PDF.
\makeatletter
\newcommand{\DeclareUnicodeFallback}[2]{%
  \begingroup\lccode`\~=#1\relax
  \lowercase{\endgroup\protected\def~}{#2}%
  \catcode#1=\active}
\makeatother

\DeclareUnicodeFallback{"2194}{\ensuremath{\leftrightarrow}}  % ↔
\DeclareUnicodeFallback{"2260}{\ensuremath{\neq}}             % ≠
\DeclareUnicodeFallback{"2265}{\ensuremath{\geq}}             % ≥
\DeclareUnicodeFallback{"26A0}{\textbf{!}}                    % ⚠

% U+FE0F only asks a renderer to draw the preceding character in colour. On
% paper it carries nothing.
\DeclareUnicodeFallback{"FE0F}{}

% No font in the image has emoji, and a printed manual does not need the
% decoration. Drop the character and the space behind it, so a heading written
% "## 📚 Tutorials" sets as "Tutorials" rather than " Tutorials".
\DeclareUnicodeFallback{"1F3D7}{\ignorespaces}  % 🏗
\DeclareUnicodeFallback{"1F4A1}{\ignorespaces}  % 💡
\DeclareUnicodeFallback{"1F4D6}{\ignorespaces}  % 📖
\DeclareUnicodeFallback{"1F4DA}{\ignorespaces}  % 📚
\DeclareUnicodeFallback{"1F527}{\ignorespaces}  % 🔧
\DeclareUnicodeFallback{"1F916}{\ignorespaces}  % 🤖
